\documentclass[11pt]{article}
\usepackage{explorer}
\newcommand{\ExplorerLanguage}{EN}
\newcommand{\ExplorerNumber}{12}
\newcommand{\ExplorerKind}{Progressive exercises}
\begin{document}
\ExplorerTitle{Upper high school · analysis and algorithms}{Bracketing a root by bisection}
\ExplorerSection{Given information}
For exercise 1: seek a zero of \(f(x)=x^{2}-2\) on [1,2]. After two halvings, the retained interval is [1.25,1.5].\par
\vspace{5pt}
\ExplorerSection{Exercise 1}
Carry out a third step of the example.\par
\vspace{5pt}
\ExplorerSection{Exercise 2}
Starting with width 1, how many halvings guarantee width at most \(\frac{1}{1000}\)?\par
\vspace{5pt}
\ExplorerSection{Exercise 3}
Why does the sign change of \(f(x)=\frac{1}{x}\) between −1 and 1 not prove the existence of a zero?\par
\vspace{5pt}
\ExplorerSection{Exercise 4}
A root lies in [1.4,1.5]. What maximum error is guaranteed by midpoint 1.45?\par
\vspace{5pt}
\end{document}
